\subsection{Attempt HS1: the direction of an ordinary Eisenstein-ideal bound}

A possible all-weight route is to replace finite Hecke determinants
by an ordinary Eisenstein ideal and a special-value formula. The
specific result checked here does not supply the needed unit
exclusion.

Hsieh, \emph{Eisenstein congruence on unitary groups and Iwasawa
main conjectures for CM fields}, JAMS 27 (2014), 753--862,
Theorem~7.18, Definition~7.19 and Corollary~7.20
(author manuscript pp.~91--92), constructs an ordinary cuspidal
congruence from a nonunit product
\[
 L^S(\Psi)=L_p^S(-1,\Psi_+\tau_{K/F})L_p^S(\Psi,\Sigma).
\]
Besides primitivity away from $p$, the theorem imposes the
stated conductor, ordinary CM-type and $\mu$-invariant conditions.
The construction and its localized congruence ideal explicitly
assume that this product is not a unit. The corollary gives
\[
                \operatorname{Eis}(\Psi,S)\subseteq(L^S(\Psi)).
\]
The introduction, inequalities (L$\mid$E) and (E$\mid$S),
and Theorem~3 give the corresponding lower bounds on
congruence valuations and Selmer lengths. These exact
statements, and the complete short proof of Theorem~7.18,
were read; no proof audit of the entire article is claimed.

The direction matters. In a discrete valuation ring,
$I\subseteq(L)$ says $v(I)\geq v(L)$. If $L$ is a unit,
that inequality alone places no restriction on $I$:
$I=(\ell)\subseteq(1)$ is a concrete counterexample to
the proposed inference. An opposite inclusion, an equality,
or a separately proved upper bound on the appropriate
congruence module would be needed to infer a unit ideal
from a unit special value. Moreover, the actual theorem's
localized ideal is constructed in the nonunit case, so
its setup cannot simply be reused unchanged in the unit case.

No such opposite inclusion or equality for the precise
finite-level compact-unitary comparison of this report
has been established. Nor has the required specialization
from Hsieh's CM Selmer group to the rational strict negative
Tate Selmer group been proved. This approach therefore leaves
an actual arithmetic comparison to establish; it does not
turn the computed generalized-Bernoulli units into an
all-weight vanishing theorem.

The primary author-hosted manuscript is
\url{https://www.math.ntu.edu.tw/~mlhsieh/research/IMCCM.pdf};
its retrieval hash and local source audit are retained in the
evidence. The publication is cited as an imported theorem.

\paragraph{Attempt NV: simultaneous units and specialization.}
Ono--Skinner, \emph{Fourier coefficients of half-integral weight
modular forms modulo $\ell$}, Ann.\ of Math.\ 147 (1998),
453--470, main Theorem (author manuscript p.~2), treats one
good eigenform and all but finitely many primes. Its local
Corollary~2 (manuscript p.~3) still treats one form. These
statements do not give the same-discriminant product (FV.8),
nor a prime bound uniform in the varying weight. The primary
manuscript is
\url{https://uva.theopenscholar.com/files/ken-ono/files/029_8.pdf}.
Only the exact statements and their stated scopes are used here;
no new full-proof audit is claimed.

Castella--Liu--Wan's six-page \emph{Addendum to: Iwasawa--Greenberg
main conjecture for nonordinary modular forms and Eisenstein
congruences on GU(3,1)}, \S1, pp.~1--2, explains that the corrected
tower nonvanishing input gives infinitely many twists, insufficient
for their previous three-value selection. They use two auxiliary
families to recover a height-one assertion; \S5, Theorem~5.0.1,
p.~6 records the repaired conclusion. The primary source is
\url{https://web.math.ucsb.edu/~castella/CLW-addendum.pdf}.
The affected Hsieh special-value paper there is the 2014
\emph{Documenta} paper; it is distinct from the JAMS paper
cited in Attempt HS1. We do not import an uncorrected
almost-all-twists assertion.

This suggests a two-family replacement for simultaneous units,
but KT-FIELD-SPECIALIZATION-LIMIT proves that height-one avoidance
alone cannot settle every specialized integral group here.
The needed additional argument must control integral specialization,
including higher-codimension contributions; merely substituting
two independent variables for (FV.8) does not prove it.
